\chapter{Interview Question Bank}
\index{interview questions}\index{question bank}%
This appendix consolidates the Interview Q\&A sections of Chapters~1--24 into one bank, grouped
by domain and tagged by difficulty: \textbf{(F)} foundational, \textbf{(I)} intermediate,
\textbf{(A)} advanced. Answers are abbreviated to the key points an interviewer listens for; the
chapter reference points to the full discussion. Practice out loud---a correct answer you cannot
say in ninety seconds is not yet an interview answer.

\section{Storage and the Data Lake}
\begin{enumerate}
    \item \textbf{(F)} \textbf{Why is S3 a better data-lake foundation than cluster-attached disks?}
    Decouples durable storage from ephemeral compute; many engines read the same data; lifecycle,
    replication, encryption, and policy are independent of any cluster. (Ch.~1)
    \item \textbf{(F)} \textbf{When would you avoid One Zone-IA?}
    For the only copy of critical or regulated data---single-AZ storage suits re-creatable
    derivatives and secondary copies. (Ch.~1)
    \item \textbf{(I)} \textbf{What does S3 strong consistency solve---and not solve?}
    Immediate visibility of completed writes and lists. It is not multi-object transactions and
    does not repair partial job output. (Ch.~1)
    \item \textbf{(I)} \textbf{Object Lock: governance or compliance mode?}
    Governance when a documented administrative bypass exists; compliance for certain, immutable
    WORM retention. (Ch.~1)
    \item \textbf{(I)} \textbf{DataSync or Snowball for 10 PB?}
    Snow Family for the initial bulk transfer; DataSync for ongoing deltas. (Ch.~1)
\end{enumerate}

\section{Databases and Caching}
\begin{enumerate}
    \item \textbf{(F)} \textbf{RDS Multi-AZ versus a Read Replica?}
    Multi-AZ: synchronous standby for HA/failover. Read replica: asynchronous, for read scaling
    and migration staging; can lag. (Ch.~2)
    \item \textbf{(I)} \textbf{Why can Aurora fail over faster than traditional databases?}
    Compute separates from a shared, replicated cluster volume---promotion needs no volume copy.
    (Ch.~2)
    \item \textbf{(I)} \textbf{When is Aurora Global Database justified?}
    Low-latency cross-region reads or regional DR beyond backups; still needs application failover
    planning. (Ch.~2)
    \item \textbf{(I)} \textbf{What does Aurora Serverless v2 not solve?}
    Schema tuning, connection management, transaction design, maximum-capacity planning.
    (Ch.~2)
    \item \textbf{(I)} \textbf{SCT's role in an Oracle$\rightarrow$Aurora PostgreSQL migration?}
    Assess and convert schema; flag manual remediation for incompatible SQL, procedures, types.
    (Ch.~2)
    \item \textbf{(A)} \textbf{Why does CDC enable near-zero-downtime migration?}
    The target stays current while the source serves traffic; cutover waits only for lag and
    validation thresholds. (Ch.~2)
    \item \textbf{(F)} \textbf{How is DynamoDB different from a relational database?}
    Managed key-value/document store modeled around known access patterns; no joins or ad hoc
    relational queries. (Ch.~3)
    \item \textbf{(F)} \textbf{What makes a good DynamoDB partition key?}
    High cardinality, even traffic distribution, groups items queried together; avoids hot keys.
    (Ch.~3)
    \item \textbf{(F)} \textbf{When choose on-demand capacity?}
    Unknown, bursty, early-stage traffic; provisioned + autoscaling for predictable steady load.
    (Ch.~3)
    \item \textbf{(I)} \textbf{LSI versus GSI?}
    LSI shares the partition key, alternate sort key, created with the table. GSI has its own key
    schema, can be added later, supports new access patterns. (Ch.~3)
    \item \textbf{(I)} \textbf{Why are sparse GSIs useful?}
    Only items with the index key appear---efficient queues, status views, entity-type slices.
    (Ch.~3)
    \item \textbf{(F)} \textbf{What are DynamoDB Streams for?}
    Item-level change feed: cache invalidation, materialized views, search indexing, audit, lake
    ingestion. (Ch.~3)
    \item \textbf{(I)} \textbf{DynamoDB TTL caveat?}
    Expiration is asynchronous---fine for cleanup, wrong for exact-time deletion deadlines.
    (Ch.~3)
    \item \textbf{(A)} \textbf{What extra design do Global Tables require?}
    Conflict-aware write rules, regional observability, idempotent consumers, multi-region-aware
    routing. (Ch.~3)
    \item \textbf{(F)} \textbf{What problem does ElastiCache solve?}
    Serves repeated reads from memory so databases/APIs see fewer requests at lower latency.
    (Ch.~4)
    \item \textbf{(I)} \textbf{Redis-compatible ElastiCache versus Memcached?}
    Redis: rich structures, replication, clustering, persistence. Memcached: simple ephemeral
    key-value caching. (Ch.~4)
    \item \textbf{(F)} \textbf{What is cache-aside?}
    Check cache, on miss load from the database, store with TTL; the database remains the source
    of truth. (Ch.~4)
    \item \textbf{(A)} \textbf{Why is cache invalidation hard?}
    Must track every write affecting a key and race readers, writers, retries, and TTL expiry.
    (Ch.~4)
    \item \textbf{(I)} \textbf{DAX versus ElastiCache?}
    DAX: DynamoDB-specific, API-compatible read cache. ElastiCache: general Redis/Valkey/Memcached
    for many sources. (Ch.~4)
    \item \textbf{(I)} \textbf{Can DAX serve strongly consistent reads?}
    No---strong reads bypass DAX to DynamoDB; design read-after-write paths explicitly. (Ch.~4)
    \item \textbf{(I)} \textbf{AppSync versus API Gateway?}
    AppSync: GraphQL with typed schemas and subscriptions. API Gateway: REST/HTTP/webhook
    ingestion. (Ch.~4)
    \item \textbf{(I)} \textbf{Which metrics show cache health?}
    Hit ratio, latency, evictions, memory pressure, CPU, connections, replication lag, error rate,
    downstream database load. (Ch.~4)
\end{enumerate}

\section{Warehousing and Lakehouse Analytics}
\begin{enumerate}
    \item \textbf{(I)} \textbf{Why declare the grain before adding fact columns?}
    Every measure must be true at exactly one row of the fact; mixed grains double-count. (Ch.~5)
    \item \textbf{(I)} \textbf{Periodic versus accumulating snapshot facts?}
    Periodic: state at intervals, never updated. Accumulating: one workflow instance, milestone
    columns updated as it advances. (Ch.~5)
    \item \textbf{(F)} \textbf{When is SCD Type 2 required over Type 1?}
    When the question depends on the attribute value at event time; pure corrections use Type 1.
    (Ch.~5)
    \item \textbf{(I)} \textbf{COPY versus Auto-Copy for micro-batch loads?}
    Auto-Copy monitors a prefix and loads on arrival; manual COPY for controlled, manifest-driven
    bulk loads. (Ch.~5)
    \item \textbf{(I)} \textbf{When does Redshift Serverless beat provisioned RA3 on cost?}
    Spiky/intermittent utilization where idle hours dominate; steady 24/7 amortizes provisioned
    better. (Ch.~5)
    \item \textbf{(A)} \textbf{What does the leader node do during COPY?}
    Plans and assigns input files to compute nodes, which ingest into their slices in parallel;
    stores no user data. (Ch.~5)
    \item \textbf{(I)} \textbf{Which distribution style avoids data movement on large joins?}
    KEY on the join column for both tables---collocated join, nothing crosses the interconnect.
    (Ch.~6)
    \item \textbf{(A)} \textbf{Compound versus interleaved sort keys?}
    Interleaved for varied filter-column combinations; compound for prefix-ordered (typically
    date-first) access. (Ch.~6)
    \item \textbf{(I)} \textbf{What does concurrency scaling solve?}
    Bursty read demand (dashboard refresh storms) with transient read-only clusters, billed per
    second. (Ch.~6)
    \item \textbf{(F)} \textbf{When does a materialized view beat a plain view?}
    Repeated expensive aggregations tolerating a refresh cadence---classic dashboards. (Ch.~6)
    \item \textbf{(A)} \textbf{Rising WLM queue wait time tells you to change what?}
    Capacity: concurrency scaling, service-class memory, cheaper queries, or throttled
    demand---wait time is demand exceeding supply, not a query bug. (Ch.~6)
    \item \textbf{(I)} \textbf{Why is ALL distribution dangerous for large tables?}
    Full copy per node: storage multiplies, loads rewrite every copy; justified only for small
    dimensions. (Ch.~6)
    \item \textbf{(I)} \textbf{Why does Parquet lower Athena cost versus CSV?}
    Per-byte-scanned billing; columnar + compressed reads only referenced column blocks---often
    10--100$\times$ fewer bytes. (Ch.~7)
    \item \textbf{(I)} \textbf{What does a Glue crawler infer, and when does it infer wrong?}
    Columns, types, partitions by sampling. Fails on mixed schemas in one prefix or
    unrepresentative samples. (Ch.~7)
    \item \textbf{(I)} \textbf{Spectrum versus Athena over the same S3 data?}
    Spectrum when joining lake with warehouse tables in one statement; Athena for lake-only ad hoc
    with no cluster. (Ch.~7)
    \item \textbf{(F)} \textbf{How does partition pruning cut bytes scanned?}
    Engine skips files outside the partition filter---seven day-partitions instead of five years.
    (Ch.~7)
    \item \textbf{(A)} \textbf{What makes Iceberg a lakehouse format beyond Parquet files?}
    Metadata layer: atomic commits, schema/partition evolution, hidden partitioning, snapshot time
    travel. (Ch.~7)
    \item \textbf{(I)} \textbf{Why crawl the curated zone, not the raw zone?}
    Inference assumes uniform layout; curated is compacted/typed/partitioned. Raw deserves
    hand-written DDL. (Ch.~7)
    \item \textbf{(I)} \textbf{What does \texttt{IAMAllowedPrincipals} do, and why revoke it?}
    Makes new catalog objects IAM-governed, bypassing Lake Formation; until revoked, LF filters do
    not bind. (Ch.~8)
    \item \textbf{(I)} \textbf{LF-Tags versus named-resource grants?}
    Tags bind policy to expressions---every correctly tagged table inherits it; names scale with
    table count, tags with data classes. (Ch.~8)
    \item \textbf{(I)} \textbf{What does column-level security protect that table grants cannot?}
    Sensitive attributes on shareable tables---segments visible, \texttt{pii\_email} absent, no
    view or copy. (Ch.~8)
    \item \textbf{(A)} \textbf{How does Redshift data sharing avoid copies?}
    Consumers query producer managed storage live and read-only---no ETL, no duplicate storage, no
    drift. (Ch.~8)
    \item \textbf{(I)} \textbf{Why is a clean room safer than exchanging CSV exports?}
    No party receives raw records; approved aggregates only, minimum group sizes, full query logs.
    (Ch.~8)
    \item \textbf{(A)} \textbf{Why must both IAM and Lake Formation allow a lake query?}
    Grants are additive gates: IAM governs API calls, LF governs catalog objects; effective
    permission is the intersection. (Ch.~8)
    \item \textbf{(F)} \textbf{What does SPICE accelerate, and its limits?}
    Dashboard dataset queries via in-memory snapshots; bounded by SPICE allocation and refresh
    schedule. (Ch.~16)
    \item \textbf{(I)} \textbf{How does QuickSight row-level security work?}
    A rules dataset maps users/groups to allowed values, applied at query time---one dataset, many
    audiences. (Ch.~16)
    \item \textbf{(F)} \textbf{What is reverse ETL and what does it feed?}
    Syncs curated warehouse results (segments, scores) back into CRMs, marketing, support systems.
    (Ch.~16)
    \item \textbf{(A)} \textbf{Why must identity resolution precede CRM audience syncs?}
    Warehouse keys must map to the right person; uncontrolled identifiers create mis-targeted
    outreach. (Ch.~16)
    \item \textbf{(I)} \textbf{Hourly versus daily reverse-ETL sync?}
    Weigh API rate limits and cost against freshness value; engineer to the loosest cadence that
    meets the stated requirement. (Ch.~16)
    \item \textbf{(I)} \textbf{When is direct query justified over SPICE?}
    Freshness in minutes and a source that bears per-interaction load---operational monitors, not
    company-wide dashboards. (Ch.~16)
\end{enumerate}

\section{ETL, CDC, and Orchestration}
\begin{enumerate}
    \item \textbf{(I)} \textbf{Core versus task nodes: which can run on Spot?}
    Task nodes---no HDFS data, interruption only reschedules tasks. Core/master on Spot risks job
    or cluster failure. (Ch.~9)
    \item \textbf{(F)} \textbf{How does EMRFS decouple storage from compute?}
    S3 as the cluster file system; clusters become disposable compute over durable objects. (Ch.~9)
    \item \textbf{(I)} \textbf{When EMR on EKS over EMR on EC2?}
    When Kubernetes is the org's compute standard and Spark should share its scheduling, quotas,
    networking. (Ch.~9)
    \item \textbf{(F)} \textbf{What are bootstrap actions for?}
    Setup scripts on every node at provisioning---reproducible transient clusters from code.
    (Ch.~9)
    \item \textbf{(I)} \textbf{Why prefer transient clusters for recurring batch?}
    Cluster-hours track job-hours; a 2-hour daily job bills 2 node-hours/day, not 24. (Ch.~9)
    \item \textbf{(I)} \textbf{What stays on local HDFS in a well-designed EMR job?}
    Job-local scratch (shuffle, spill); all durable inputs/outputs live in S3. (Ch.~9)
    \item \textbf{(A)} \textbf{What does a DynamicFrame offer over a DataFrame?}
    Self-describing records with \texttt{resolveChoice}/\texttt{applyMapping}/
    \texttt{relationalize} for messy, evolving schemas and catalog-aware sinks. (Ch.~10)
    \item \textbf{(I)} \textbf{What do job bookmarks guarantee?}
    Inputs already processed are not reprocessed. They do not make the sink idempotent. (Ch.~10)
    \item \textbf{(I)} \textbf{When is Flex execution cheaper?}
    Non-urgent jobs (backfills, slack-heavy overnight batch) that tolerate delayed starts. (Ch.~10)
    \item \textbf{(F)} \textbf{Why Glue over EMR for analyst-built pipelines?}
    No clusters, visual designer generating real code, catalog integration, per-second billing.
    (Ch.~10)
    \item \textbf{(F)} \textbf{What compute is one DPU?}
    4 vCPU + 16 GB, billed per second with a one-minute minimum. (Ch.~10)
    \item \textbf{(I)} \textbf{DataBrew's best use?}
    Interactive profiling and recipe prototyping; also a schedulable quality gate inside a
    pipeline. (Ch.~10)
    \item \textbf{(I)} \textbf{Which CloudWatch metrics reveal CDC lag on a ``running'' task?}
    \texttt{CDCLatencySource} and \texttt{CDCLatencyTarget}; task status says nothing about
    freshness. (Ch.~11)
    \item \textbf{(A)} \textbf{BACKWARD versus FULL compatibility: which permits field deletion?}
    BACKWARD (new reads old). FULL forbids deletions old consumers cannot tolerate. (Ch.~11)
    \item \textbf{(I)} \textbf{Who owns a data contract?}
    Producer owns and versions the schema; consumers sign off; the registry enforces mechanics,
    the workflow enforces intent. (Ch.~11)
    \item \textbf{(I)} \textbf{Why enforce schema at ingestion, not query time?}
    By query time bad data is durable---and possibly replicated and transformed downstream.
    (Ch.~11)
    \item \textbf{(F)} \textbf{One breaking and one non-breaking schema change?}
    Breaking: delete/rename/retype a field. Non-breaking: add an optional field with a default.
    (Ch.~11)
    \item \textbf{(A)} \textbf{Why stage DMS output in S3 before Redshift?}
    Replayable change log decouples replication from warehouse availability; MERGEs re-runnable
    and auditable. (Ch.~11)
    \item \textbf{(I)} \textbf{What does \texttt{catchup=False} prevent?}
    Backfill storms---every missed historical run launching on deploy against production. (Ch.~12)
    \item \textbf{(I)} \textbf{When are Step Functions better than Airflow?}
    Event-driven, high-rate workflows of AWS service calls where per-transition billing beats an
    always-on scheduler. (Ch.~12)
    \item \textbf{(F)} \textbf{What do retries with backoff protect against?}
    Transient failures (throttling, brief hiccups). Deterministic failures need the failure path.
    (Ch.~12)
    \item \textbf{(I)} \textbf{Why tune a sensor's \texttt{poke\_interval}?}
    Each poll is billed LIST requests; match the interval to the expected arrival cadence. (Ch.~12)
    \item \textbf{(I)} \textbf{When should data move through S3 instead of XComs?}
    Always, for datasets; XComs (and the 256 KB state-machine payload) suit small signals only.
    (Ch.~12)
    \item \textbf{(A)} \textbf{Why does the alert task need \texttt{trigger\_rule="one\_failed"}?}
    Default rules skip the alert on upstream skip---exactly the runs it exists for. (Ch.~12)
\end{enumerate}

\section{Streaming}
\begin{enumerate}
    \item \textbf{(I)} \textbf{Which two throughput limits size a Kinesis stream?}
    Per-shard 1 MB/s ingress \emph{and} 1{,}000 records/s; shard count is the max plus headroom.
    (Ch.~13)
    \item \textbf{(I)} \textbf{Why does partition-key choice control per-entity ordering?}
    The key hashes to a shard; a shard is the ordered unit. Skewed keys concentrate load. (Ch.~13)
    \item \textbf{(I)} \textbf{What does rising \texttt{IteratorAgeMilliseconds} warn of?}
    Consumers falling behind finite retention---alarm before the retention horizon. (Ch.~13)
    \item \textbf{(A)} \textbf{How do you build exactly-once effects on at-least-once delivery?}
    Idempotent consumers: dedup on event ID or upsert by key; poison records to a DLQ. (Ch.~13)
    \item \textbf{(I)} \textbf{Kinesis versus MSK billing models?}
    Kinesis: shard-hours + record volume (capacity). MSK: broker instance-hours + storage (a
    fleet). (Ch.~13)
    \item \textbf{(I)} \textbf{When does SQS FIFO suffice and Kinesis is overkill?}
    One consumer per message, no replay, throughput within FIFO quotas. (Ch.~13)
    \item \textbf{(F)} \textbf{Name three native Firehose destinations.}
    S3, Redshift (via intermediate S3 COPY), OpenSearch---plus Splunk and HTTP endpoints. (Ch.~14)
    \item \textbf{(I)} \textbf{Why enable source-record backup on a delivery stream?}
    Pre-transformation originals preserved regardless of delivery success: replay, audit,
    recovery. (Ch.~14)
    \item \textbf{(I)} \textbf{What do inline Lambda transformations add?}
    Execution latency per batch plus Lambda GB-seconds and the per-GB surcharge; per-record network
    calls explode both. (Ch.~14)
    \item \textbf{(I)} \textbf{Why convert JSON to Parquet during Firehose delivery?}
    No separate conversion job; the lake is columnar at landing, so scanned-bytes savings apply
    from the first query. (Ch.~14)
    \item \textbf{(I)} \textbf{Where does data land when Redshift delivery keeps failing?}
    The intermediate S3 staging prefix---needs lifecycle rules and a
    \texttt{DeliveryToRedshift.Success} alarm. (Ch.~14)
    \item \textbf{(A)} \textbf{What can a Firehose Lambda transform never correctly do?}
    Cross-record state: dedup windows, joins, aggregations---stream-processor work. (Ch.~14)
    \item \textbf{(A)} \textbf{Why does unbounded keyed state kill a streaming job?}
    Checkpoints grow until they cannot complete; restores take hours; RocksDB spills. Every
    stateful operator needs a bound. (Ch.~15)
    \item \textbf{(I)} \textbf{Checkpoint versus savepoint before redeployment?}
    Savepoint---manual, stable, portable across graph changes; checkpoints serve automatic
    recovery. (Ch.~15)
    \item \textbf{(A)} \textbf{What bounds join-state retention in a stream--stream interval join?}
    The interval predicate plus watermarks on both inputs; past the watermark's reach, state is
    dropped. (Ch.~15)
    \item \textbf{(A)} \textbf{Which two components give end-to-end exactly-once to Kafka or S3?}
    Checkpointing plus a two-phase-commit sink---or an idempotent-by-key sink achieving the same
    effect. (Ch.~15)
    \item \textbf{(I)} \textbf{When prefer at-least-once plus downstream dedup?}
    Non-transactional sinks, throughput over strictness, or a natural dedup key---simpler, faster,
    honest. (Ch.~15)
    \item \textbf{(I)} \textbf{What does a watermark assert?}
    ``No event older than this timestamp will arrive''---a heuristic trading result latency against
    late-data tolerance. (Ch.~15)
\end{enumerate}

\section{Machine Learning and MLOps}
\begin{enumerate}
    \item \textbf{(I)} \textbf{What does Redshift ML delegate to SageMaker?}
    Everything past export: unload to S3, Autopilot training/selection, winner compiled back into a
    SQL function. (Ch.~17)
    \item \textbf{(F)} \textbf{Why is in-database ML attractive for SQL-only pipelines?}
    Features, training, scoring live where the data is---no extraction pipelines, no endpoints, no
    new operational surface. (Ch.~17)
    \item \textbf{(I)} \textbf{What does \texttt{CREATE MODEL}'s \texttt{FUNCTION} clause define?}
    The scoring function's name and signature, mapping arguments to the training query's feature
    columns. (Ch.~17)
    \item \textbf{(I)} \textbf{Why F1 over accuracy for imbalanced churn?}
    5\% positives: always-negative scores 95\% while detecting nothing; F1 balances precision and
    recall on the minority class. (Ch.~17)
    \item \textbf{(A)} \textbf{When graduate from Redshift ML to full SageMaker?}
    Custom architectures, fine hyperparameter control, formal explainability, low-latency
    real-time inference. (Ch.~17)
    \item \textbf{(A)} \textbf{What is target leakage and how do you catch it?}
    A feature unknowable at prediction time in training; review the training query column-by-column
    against ``knowable before the outcome?'' (Ch.~17)
    \item \textbf{(I)} \textbf{Real-time vs.\ serverless vs.\ Batch Transform for sporadic traffic?}
    Serverless: scale to zero, pay per request, accept cold starts. Real-time for steady SLAs;
    Batch Transform for bulk offline. (Ch.~18)
    \item \textbf{(F)} \textbf{What does Data Wrangler export?}
    A repeatable artifact---Processing Job script, pipeline step, or notebook---not a one-off
    session. (Ch.~18)
    \item \textbf{(I)} \textbf{When do you need a custom container over built-in XGBoost?}
    Framework/architecture/dependencies with no managed equivalent; otherwise stay on maintained
    images. (Ch.~18)
    \item \textbf{(I)} \textbf{How does Autopilot pick its final model?}
    Explores preprocessors, algorithms, hyperparameters; evaluates on validation splits against
    the objective; ranks a winner with inspectable candidate notebooks. (Ch.~18)
    \item \textbf{(A)} \textbf{Why route training-data access through Lake Formation?}
    Column masks apply to extraction; PII never reaches staging; every training set is governed
    and reproducible. (Ch.~18)
    \item \textbf{(I)} \textbf{What is managed Spot training and its code requirement?}
    Spare capacity at a steep discount with interruption---requires checkpointing to S3 to resume
    rather than restart. (Ch.~18)
    \item \textbf{(I)} \textbf{Online versus offline Feature Store?}
    Offline (S3, point-in-time joins) builds training sets; online (millisecond
    \texttt{GetRecord}) serves inference; one definition feeds both---the skew defense. (Ch.~19)
    \item \textbf{(F)} \textbf{What does the Model Registry gate?}
    Approval: deployment keys on the \texttt{Approved} state of a versioned, metrics-carrying
    package. (Ch.~19)
    \item \textbf{(A)} \textbf{Data drift versus concept drift?}
    Data: input distributions change (features alone detect). Concept: input$\rightarrow$outcome
    relationship changes (needs delayed ground truth). (Ch.~19)
    \item \textbf{(A)} \textbf{Why gate automatic retraining behind human approval?}
    Drift may mean the world changed \emph{or} ingestion broke; auto-retraining on corrupt data
    automates worse models. (Ch.~19)
    \item \textbf{(I)} \textbf{What baseline does Model Monitor need?}
    Per-feature statistics and constraints from training data, refreshed on a retraining cadence.
    (Ch.~19)
    \item \textbf{(I)} \textbf{What does pipeline step caching buy?}
    Reruns with unchanged inputs reuse outputs---cheaper iteration with input-fingerprint
    correctness. (Ch.~19)
    \item \textbf{(F)} \textbf{Textract versus Comprehend for scanned-form key--value pairs?}
    Textract \texttt{AnalyzeDocument FORMS}; Comprehend works on already-digital text downstream.
    (Ch.~20)
    \item \textbf{(F)} \textbf{Why trigger AI APIs from S3 events via Lambda?}
    Arrival-driven, no polling, per-object scale-out matching spiky uploads. (Ch.~20)
    \item \textbf{(I)} \textbf{When switch from sync to async Textract?}
    Multi-page PDFs/TIFFs; async signals completion via SNS. (Ch.~20)
    \item \textbf{(I)} \textbf{How make extracted document metadata searchable?}
    Persist normalized JSON, index into OpenSearch, register results in the catalog for SQL.
    (Ch.~20)
    \item \textbf{(A)} \textbf{PII risks of persisting raw OCR output?}
    OCR is faithful---IDs and addresses land in the JSON; restricted zone, field-level redaction,
    at-rest discovery scanning. (Ch.~20)
    \item \textbf{(I)} \textbf{How keep retried pipelines from re-billing API calls?}
    Idempotency keyed on object identity (etag/version); skip already-processed inputs. (Ch.~20)
\end{enumerate}

\section{Security, Governance, and Compliance}
\begin{enumerate}
    \item \textbf{(I)} \textbf{Gateway versus interface VPC endpoint for S3?}
    Gateway: free route-table entry for S3/DynamoDB. Interface (PrivateLink): other services, billed
    per AZ-hour plus data. (Ch.~21)
    \item \textbf{(I)} \textbf{What does a permissions boundary limit?}
    A role's maximum permissions regardless of identity policies---the intersection is effective;
    makes delegated role creation safe. (Ch.~21)
    \item \textbf{(F)} \textbf{Why is automatic Secrets Manager rotation better than static keys?}
    A leaked credential's value is bounded to the rotation window; static keys leak forever.
    (Ch.~21)
    \item \textbf{(I)} \textbf{How do SCPs differ from IAM policies?}
    SCPs cap what any account---including its admins---may do; IAM policies grant within an
    account. (Ch.~21)
    \item \textbf{(A)} \textbf{What exfiltration path does PrivateLink close that security groups
    cannot?}
    Endpoint policies restrict the \emph{service destinations} themselves---your network cannot
    write to a stranger's bucket. (Ch.~21)
    \item \textbf{(A)} \textbf{Why does KMS complete the perimeter?}
    Key policies restrict decryption to org principals; copied data arrives as unusable ciphertext.
    (Ch.~21)
    \item \textbf{(I)} \textbf{How do LF-Tags simplify cross-account sharing?}
    Grants bind to tag expressions; correctly tagged shared tables inherit policy in every account.
    (Ch.~22)
    \item \textbf{(I)} \textbf{Why consult lineage when a quality check fails?}
    Names what consumed the bad batch and which upstream change introduced it---a directed fix, not
    a two-day hunt. (Ch.~22)
    \item \textbf{(I)} \textbf{\texttt{ColumnExists} versus \texttt{Completeness}?}
    Structure versus content; a production ruleset needs both families. (Ch.~22)
    \item \textbf{(F)} \textbf{In a data mesh, who owns a data product's quality?}
    The publishing domain; the center provides platform, taxonomy, audit. (Ch.~22)
    \item \textbf{(I)} \textbf{What does a DQ rule's threshold clause buy?}
    Distinguishes a mild dip (report) from a cliff (fail)---monitoring without halting on noise.
    (Ch.~22)
    \item \textbf{(A)} \textbf{Why enforce contracts at registration, not query time?}
    Registration declares producer intent; query time means violations are already durable bad
    data. (Ch.~22)
    \item \textbf{(I)} \textbf{What does Macie detect beyond regex patterns?}
    Checksums (Luhn), proximity keywords, format validation; also bucket posture inventory.
    (Ch.~23)
    \item \textbf{(I)} \textbf{Why quarantine suspect files with Object Lock instead of deleting?}
    Preserves immutable evidence for investigation and audit; deletion can violate retention.
    (Ch.~23)
    \item \textbf{(A)} \textbf{Tokenization versus hashing: which is reversible?}
    Tokenization, via the controlled vault. Low-entropy hashes are effectively reversible to
    attackers---use real tokenization when re-identification is required. (Ch.~23)
    \item \textbf{(I)} \textbf{What do CloudTrail data events record for S3?}
    Object-level API activity per principal and object---the lake's access ledger, delivered to an
    independent log-archive account. (Ch.~23)
    \item \textbf{(I)} \textbf{Why tune Macie custom identifiers before rollout?}
    False-positive economics: over-broad identifiers train teams to ignore alerts. Tune on samples;
    review precision monthly. (Ch.~23)
    \item \textbf{(A)} \textbf{Why is crypto-shredding superior to enumerated deletion?}
    One key deletion renders every copy---versions, replicas, snapshots, archives---unrecoverable in
    a single auditable act. (Ch.~23)
\end{enumerate}

\section{DataOps and Delivery}
\begin{enumerate}
    \item \textbf{(I)} \textbf{Why store Terraform state remotely with locking?}
    Shared state is the source of truth; locking prevents divergent applies; local state is lost
    with laptops and leaks attribute values. (Ch.~24)
    \item \textbf{(I)} \textbf{What do dbt tests validate that schemas cannot?}
    Semantics: uniqueness, referential integrity, accepted ranges, freshness. Schema guarantees
    shape; tests guarantee the data is right. (Ch.~24)
    \item \textbf{(I)} \textbf{Why promote dbt runs Dev$\rightarrow$Prod instead of one shared
    workspace?}
    Every change is rehearsed against a real environment; one shared workspace makes production
    the rehearsal. (Ch.~24)
    \item \textbf{(A)} \textbf{Why prefer OIDC federation over long-lived AWS keys in CI?}
    Short-lived, audience-scoped credentials per run---nothing stored to leak, nothing to rotate.
    (Ch.~24)
    \item \textbf{(I)} \textbf{What must the PR-check workflow prove before merge?}
    Reviewable \texttt{terraform plan}, \texttt{dbt build} passing against CI data, policy checks
    green---\texttt{main} is always deployable. (Ch.~24)
    \item \textbf{(A)} \textbf{Summarize expand-contract in one breath.}
    Add new alongside old; dual-write and backfill while reads stay old; switch reads behind a
    shim; verify consumers; drop the old column in a later release. (Ch.~24)
\end{enumerate}

\begin{tipbox}
Interviewers score the \emph{reasoning}, not the product name. Answer in the pattern this book
teaches: state the constraint, name the trade-off, then choose---and mention what you would
monitor to know the choice was wrong.
\end{tipbox}
